\documentclass[11pt]{article}
\usepackage{amsmath, amssymb}
\usepackage[margin=1in]{geometry}

\title{GP Moments Reference: Before and After Conditioning}
\author{summary of what happens when evaluating $U_{DA}$ of an rescaled image }
\date{February 2026}

\begin{document}
\maketitle

%----------------------------------------------------------------------
\section{Kernel}
%----------------------------------------------------------------------

Arc-cosine kernel (order 1):
\begin{equation}\label{eq:kernel}
    K(x, y) = \frac{1}{\pi}\,\|x\|_C\;\|y\|_C \cdot J(\theta)
\end{equation}
where
\begin{align}
    \|x\|_C &= \sqrt{x^\top\! C\, x + \sigma_0^2}
        &\text{(C-norm, ``amplitude'')} \\[4pt]
    \cos\theta &= \frac{x^\top\! C\, y + \sigma_0^2}{\|x\|_C\;\|y\|_C}
        &\text{(``angle'' between $x$ and $y$)} \\[4pt]
    J(\theta) &= \sin\theta + (\pi - \theta)\cos\theta
        &\text{(angular factor, monotone decreasing on $[0,\pi]$)}
\end{align}

\textbf{Self-kernel:} Setting $y = x$ in~\eqref{eq:kernel}:
the angle is $\theta(x,x) = 0$ (since $\cos\theta = \|x\|_C^2 / \|x\|_C^2 = 1$),
and $J(0) = \sin 0 + \pi\cos 0 = \pi$, so
\begin{equation}
    K(x, x) = \frac{1}{\pi}\,\|x\|_C\;\|x\|_C \cdot \pi = \|x\|_C^2.
\end{equation}
The prior variance at any point equals its squared C-norm.

%----------------------------------------------------------------------
\section{Variational GP Posterior}
%----------------------------------------------------------------------

We have $M$ inducing points $\{\tilde z_j\}$ with variational parameters
$q(\tilde\lambda) = \mathcal{N}(m, V)$.

Define the shorthands:
\begin{align}
    \tilde K &= K(\tilde Z, \tilde Z)
        &&\text{($M\!\times\!M$ inducing-point kernel matrix)} \\
    k(x) &= \bigl[K(x, \tilde z_1),\;\ldots,\;K(x, \tilde z_M)\bigr]^\top
        &&\text{($M$-vector, cross-kernel to inducing points)} \\
    W &= \tilde K^{-1}(V - \tilde K)\,\tilde K^{-1}
        &&\text{(``correction'' matrix, fixed after training)}
\end{align}

%----------------------------------------------------------------------
\subsection{Marginal moments at $x^*$}
%----------------------------------------------------------------------

\begin{align}
    \mu(x^*) &= k(x^*)^\top \tilde K^{-1}\, m
        \label{eq:mu} \\[6pt]
    \sigma^2(x^*) &= \underbrace{K(x^*,x^*)}_{= \|x^*\|_C^2}
        + \; k(x^*)^\top\, W\, k(x^*)
        \label{eq:sigma2}
\end{align}
The first term is the prior variance; the second is the posterior
correction (negative when $V \prec \tilde K$, which shrinks the variance).

%----------------------------------------------------------------------
\subsection{Posterior covariance between $x^*$ and $x_t$}
%----------------------------------------------------------------------

\begin{equation}\label{eq:cov}
    \Sigma_q(x^*, x_t) = K(x^*, x_t) + k(x^*)^\top\, W\, k(x_t)
\end{equation}

Note: $\Sigma_q(x^*, x^*) = \sigma^2(x^*)$ and $\Sigma_q(x_t, x_t) = \sigma^2(x_t)$.

%----------------------------------------------------------------------
\subsection{Posterior correlation $\rho^2$}
%----------------------------------------------------------------------

\begin{equation}\label{eq:rho2}
    \boxed{\;
    \rho^2(x^*, x_t)
    = \frac{\Sigma_q(x^*, x_t)^2}
           {\sigma^2(x^*)\;\sigma^2(x_t)}
    \;}
\end{equation}

This is the standard Pearson correlation squared of the joint Gaussian
$[\lambda(x^*),\;\lambda(x_t)]$ under the variational posterior.
It takes values in $[0, 1]$ and measures how much knowing
$\lambda(x_t)$ tells you about $\lambda(x^*)$.

%----------------------------------------------------------------------
\section{Conditioning: Moments After Observing $\lambda(x_t)$}
%----------------------------------------------------------------------

Standard Gaussian conditioning on observing $\lambda_t$ at $x_t$:

\begin{align}
    \mu_{\mathrm{cond}}(x^*)
        &= \mu(x^*) + \frac{\Sigma_q(x^*, x_t)}{\sigma^2(x_t)}
            \bigl(\lambda_t - \mu(x_t)\bigr)
        \label{eq:mu_cond} \\[8pt]
    \sigma^2_{\mathrm{cond}}(x^*)
        &= \sigma^2(x^*) - \frac{\Sigma_q(x^*, x_t)^2}{\sigma^2(x_t)}
        \label{eq:sigma2_cond}
\end{align}

Using the definition of $\rho^2$, the conditional variance simplifies to:
\begin{equation}\label{eq:var_reduction}
    \boxed{\;
    \sigma^2_{\mathrm{cond}}(x^*)
    = \sigma^2(x^*)\;\bigl(1 - \rho^2\bigr)
    \;}
\end{equation}

So $\rho^2$ is \emph{directly} the fractional variance reduction from
conditioning. If $\rho^2 = 0.9$, conditioning removes 90\% of the
variance at $x^*$.

%----------------------------------------------------------------------
\section{Log-Firing Rate Transform}
%----------------------------------------------------------------------

The observation model is $r \sim \mathrm{Poisson}(\exp(g))$ where
$g = A\,\lambda + \lambda_0$ is the log-firing rate. The GP moments
in $g$-space are:

\begin{align}
    \mu_g &= A\,\mu + \lambda_0, &
    \sigma^2_g &= A^2\,\sigma^2
        &&\text{(marginal)} \\[4pt]
    \mu_{g,\mathrm{cond}} &= A\,\mu_{\mathrm{cond}} + \lambda_0, &
    \sigma^2_{g,\mathrm{cond}} &= A^2\,\sigma^2_{\mathrm{cond}}
        &&\text{(conditional)}
\end{align}

The variance reduction carries through:
$\sigma^2_{g,\mathrm{cond}} = \sigma^2_g\,(1 - \rho^2)$.

%----------------------------------------------------------------------
\section{DA Utility}
%----------------------------------------------------------------------

The DA utility is the drop in entropy caused by conditioning:
\begin{equation}\label{eq:uda}
    U_{\mathrm{DA}}(x^*)
    = H\!\bigl(\mu_g,\;\sigma^2_g\bigr)
    - H\!\bigl(\mu_{g,\mathrm{cond}},\;\sigma^2_{g,\mathrm{cond}}\bigr)
\end{equation}
where $H(\mu_g, \sigma^2_g) = -\sum_r p(r)\log p(r)$ is the entropy
of the Poisson--log-normal mixture (computed via Laplace approximation).

Substituting $\sigma^2_{g,\mathrm{cond}} = \sigma^2_g\,(1 - \rho^2)$
and, when $\lambda_t = \mu(x_t)$ (deterministic mode),
$\mu_{g,\mathrm{cond}} = \mu_g$:
\begin{equation}\label{eq:uda_rho}
    \boxed{\;
    U_{\mathrm{DA}}(x^*)
    = H\!\bigl(\mu_g,\;\sigma^2_g\bigr)
    - H\!\bigl(\mu_g,\;\sigma^2_g\,(1 - \rho^2)\bigr)
    \;}
\end{equation}

So the DA utility depends on exactly three quantities:
\begin{enumerate}
    \item $\mu_g = A\,\mu(x^*) + \lambda_0$ --- the operating point
        (where on the $H$ landscape) [NB: this is only fixed in this treatment cause we are dealing with the deterministic case],
    \item $\sigma^2_g = A^2\,\sigma^2(x^*)$ --- the marginal variance
        (how much entropy there is to begin with),
    \item $\rho^2$ --- how much of that variance conditioning removes
        (the horizontal shift on the landscape).
\end{enumerate}

Conditioning moves the point from $(\mu_g,\;\sigma^2_g)$ to
$(\mu_g,\;\sigma^2_g(1-\rho^2))$ --- same $\mu_g$, reduced $\sigma^2_g$.
The utility is the difference in $H$ between these two points.


%----------------------------------------------------------------------
\section{How Scaling Affects $\rho^2$}
%----------------------------------------------------------------------

Consider conditioning on $x_t$ and evaluating at $x^* = c\,x_t$
(same image, scaled by $c > 0$).

\subsection*{Exact result ($\sigma_0^2 = 0$)}

When $\sigma_0^2 = 0$, scaling preserves the angle: $\theta(c\,x_t, z) = \theta(x_t, z)$
for any $z$, because
\[
    \cos\theta(c\,x_t, z)
    = \frac{c\,x_t^\top C\, z}{\sqrt{c^2 q \cdot v_z}}
    = \frac{x_t^\top C\, z}{\sqrt{q \cdot v_z}}
    = \cos\theta(x_t, z),
    \qquad q \triangleq x_t^\top C\, x_t.
\]
Since $J(\theta)$ is unchanged and $\|c\,x_t\|_C = c\,\|x_t\|_C$,
the cross-kernel to any point scales linearly:
\begin{equation}
    K(c\,x_t, z) = c \cdot K(x_t, z)
    \qquad\Longrightarrow\qquad
    k(c\,x_t) = c \cdot k(x_t).
\end{equation}

Substituting into the posterior formulas~\eqref{eq:mu}--\eqref{eq:cov}:
\begin{align}
    \mu(c\,x_t) &= c\,\mu(x_t), \\
    \sigma^2(c\,x_t) &= c^2\,\sigma^2(x_t), \\
    \Sigma_q(c\,x_t,\; x_t) &= c\,\Sigma_q(x_t, x_t) = c\,\sigma^2(x_t).
\end{align}

Therefore
\begin{equation}
    \rho^2 = \frac{[c\,\sigma^2(x_t)]^2}{c^2\,\sigma^2(x_t) \cdot \sigma^2(x_t)} = 1.
\end{equation}

Scaling preserves the correlation \emph{exactly}, at any value of $c$. 

NOTE That the posterior moments
instead keep growing ($\mu_g \propto c$, $\sigma^2_g \propto c^2 \implies H_{marg} \to \inf$) so even if conditioning
always removes 100\% of the variance, the utility still goes to infinity.

\subsection*{With $\sigma_0^2 > 0$ (our trained model)}


The angle is no longer exactly preserved with changing $c$: $\sigma_0^2$ contributes
a constant offset to both numerator and denominator of $\cos\theta$.
The fractional error is $O(\sigma_0^2 / \|x_t\|_C^2)$ and does not
vanish with $c$.

The net effect: $\rho^2$ converges to a constant $\rho^2_\infty < 1$
as $c \to \infty$, not to~1. Conditioning removes a fixed fraction
of the variance, not all of it.

Crucially, this does \emph{not} prevent $U_{\mathrm{DA}}$ from diverging.
The moments still grow: $\mu_g \propto c$, $\sigma^2_g \propto c^2$
(with small constant fractional corrections). The conditional variance
$\sigma^2_g(1 - \rho^2_\infty)$ is a fixed fraction of $\sigma^2_g$,
so it also grows as $c^2$. Since $H$ grows roughly linearly in
$\sigma^2_g$ for large $\sigma^2_g$:
\[
    U_{\mathrm{DA}}
    = H(\mu_g,\;\sigma^2_g) - H(\mu_g,\;\sigma^2_g(1 - \rho^2_\infty))
    \;\propto\; \rho^2_\infty \cdot \sigma^2_g
    \;\propto\; c^2.
\]
The $\sigma_0^2$ offset pins the conditional point to a fixed
\emph{fraction} of $\sigma^2_g$, not to a fixed \emph{value}.
Both endpoints move outward on the $H$ landscape, and the gap
between them keeps growing.

\subsection*{For two different images}

When $x^*$ and $x_t$ are \emph{not} parallel (different natural images),
$\rho^2$ depends on their angle $\theta(x^*, x_t)$.
At large angles ($\theta \to \pi/2$), the cross-covariance
$\Sigma_q(x^*, x_t) \to 0$, so $\rho^2 \to 0$
and conditioning has no effect.
%----------------------------------------------------------------------
\section{The Normalized Kernel: What It Fixes and Where It Fails}
%----------------------------------------------------------------------

The divergence of $U_{\mathrm{DA}}$ with norm comes from the prior
variance $K(x,x) = \|x\|_C^2$ growing unbounded ( this scales $K$ and therefore the marginal moments \ref{eq:mu}, \ref{eq:sigma2}. A natural fix
is to normalize the kernel:
\begin{equation}
    \bar K(x, y) = \frac{K(x, y)}{\sqrt{K(x,x)\;K(y,y)}}
    = \frac{\frac{1}{\pi}\,\|x\|_C\,\|y\|_C\,J(\theta)}
           {\|x\|_C\,\|y\|_C}
    = \frac{J(\theta)}{\pi}.
\end{equation}

The norms cancel exactly. $\bar K$ depends only on the angle $\theta$.

\subsection*{What this fixes}

The prior variance becomes constant:
\begin{equation}
    \bar K(x, x) = \frac{J(0)}{\pi} = 1
    \qquad\text{for all $x$.}
\end{equation}
The cross-kernel entries $\bar K(x, \tilde z_j) = J(\theta_j)/\pi$
are bounded in $[0, 1]$. The posterior variance $\bar\sigma^2(x)$
is therefore bounded regardless of $\|x\|_C$.

With $\bar\sigma^2$ bounded, $\sigma^2_g = A^2\bar\sigma^2$ is
bounded, and $U_{\mathrm{DA}}$ cannot diverge by scaling.
The only way to increase utility is to improve $\rho^2$ by finding
images with small angle to $x_t$. When $\sigma_0^2 > 0$, the angle
$\theta(c\,x_t, x_t)$ is minimized at $c = 1$, so the gradient
points towards matching $x_t$ in both direction and magnitude.

\subsection*{Where it fails}

The normalized kernel discards all norm information. But image
norm carries genuine signal for neural encoding: neurons respond
differently to high-contrast and low-contrast versions of the
same pattern. By collapsing all scaled versions to the same
kernel value, $\bar K$ loses this predictive dimension.

Empirically, training with $\bar K$ on neural data drops test
correlation by $\sim$25\% compared to the unnormalized kernel.
The norm is not just an optimization nuisance --- it encodes
real structure that the GP needs for accurate predictions.


\end{document}
